Fig.~\ref{fig:sil} shows both transitions with the PX4 firmware in the loop, from the trims reached under the same commands. The front transition reaches cruise in \silFrontT{}~s and the back transition enters the hover-exit band at \silBackT{}~s, both slower than on the identified model (Section~\ref{sec:nominal}): the firmware's aircraft carries the drag of its fuselage and tail plates and a slower attitude loop than the identified one, and the policy, which is not told which plant it is flying, commands the pitch and the pusher it learned for the box, so the airspeed follows the drag of the aircraft it meets; in the back transition the last metre per second is bled by drag alone, with the pusher at zero and the rotors near hover thrust, which takes ten seconds. The angle of attack stays within $[\silFrontAlMin^\circ,\silFrontAlMax^\circ]$ and $[\silBackAlMin^\circ,\silBackAlMax^\circ]$, the flight path within $\silFrontGam^\circ$ and $\silBackGam^\circ$, and the lift fraction evaluated on the identified model never falls below \silFrontLf{} and \silBackLf{}. The pitch setpoint carries the noise of the firmware's airspeed estimate, which the attitude loop filters out of the pitch; the module certified every command (\silFrontInterv{} interventions in \silFrontUpd{} and \silBackInterv{} in \silBackUpd{} updates). The update, network and check together, took \silUpdMs{}~ms on average in the loop, on a processor shared with the parallel evaluations.
\begin{figure*}[t]\centering\includegraphics[width=0.72\textwidth]{fig11_sil.pdf}
\caption{Closed loop with the PX4 firmware in software-in-the-loop, front (left) and back (right) transitions: airspeed from the autopilot's estimator, pitch setpoint sent to the attitude loop with the pitch and angle of attack it produced, rotor and pusher commands.}\label{fig:sil}
\end{figure*}
